\documentclass[../../../include/open-logic-section]{subfiles}

\begin{document}

\olfileid{sth}{card-arithmetic}{ch}

\olsection{কন্টিনিউয়াম অনুমান}

আগের ফল থেকে ইঙ্গিত মেলে—এবং ইঙ্গিতটি ঠিক—যে অসীম
অঙ্কবাচক সংখ্যাগুলি $2$-এর ঘাতের সঙ্গে, অর্থাৎ
\olref[opps]{lem:SizePowerset2Exp} অনুযায়ী ঘাতসেট
নেওয়ার সঙ্গে, সরলভাবে সম্পর্কিত হলে অঙ্কবাচক সংখ্যার
ঘাত হিসাব \emph{সহজ} হতো। কিন্তু তারা যে সত্যিই
ঘাতসেটের সঙ্গে এমন সরলভাবে সম্পর্কিত, তা ধরে নিতে পারি না।

বিষয়টি খুলতে প্রথমে কিছু সুবিধাজনক সংকেত আনি।

\begin{defn}
$\cardsucc{\cardfont{a}}$ হলো $\cardfont{a}$-র চেয়ে
কঠোরভাবে বড় ক্ষুদ্রতম অঙ্কবাচক সংখ্যা। এর সাহায্যে
দুটি অসীম অনুক্রম সংজ্ঞায়িত করি:
\begin{align*}
	\aleph_{0} &\defis \omega &
	\beth_{0} &\defis \omega\\
	\aleph_{\alpha \ordplus 1} &\defis \cardsucc{(\aleph_{\alpha})} &
	\beth_{\alpha+1} &\defis \cardexpo{2}{\beth_{\alpha}}\\
	\aleph_{\alpha} &\defis \bigcup_{\beta< \alpha} \aleph_{\beta} &
	\beth_{\alpha} &\defis \bigcup_{\beta < \alpha}\beth_{\beta} & \text{when $\alpha$ is a limit ordinal}.
	\end{align*}
\end{defn}

$\cardsucc{\cardfont{a}}$-র সংজ্ঞাটি সার্থক। কারণ
\olref[cardinals][classing]{lem:NoLargestCardinal}
অনুযায়ী প্রত্যেক $\cardfont{a}$-র চেয়ে বড়
অঙ্কবাচক সংখ্যা আছে।
% BN-SRC-459: least greater cardinal follows from ordinal well-ordering,
% not directly from transfinite induction; recursion defines the sequences.
তাদের মধ্যে ক্ষুদ্রতমটি ক্রমসংখ্যার সুক্রমবিন্যাসে
পাওয়া যায়। অনুক্রমদুটির বাকি অংশ ট্রান্সফাইনাইট
পুনরাবৃত্তিতে সংজ্ঞায়িত হয়।

কান্টর ``$\aleph$'' সংকেতটি এনেছিলেন। এটি
\emph{আলেফ}: হিব্রু বর্ণমালার প্রথম অক্ষর, এবং
``অসীম'' অর্থের হিব্রু শব্দেরও প্রথম অক্ষর।
পিয়ার্স ``$\beth$'' সংকেতটি এনেছিলেন; \emph{বেথ}
হিব্রু বর্ণমালার দ্বিতীয় অক্ষর।\footnote{১৯০০ সালের
ডিসেম্বরে কান্টরকে লেখা একটি চিঠিতে পিয়ার্স এই সংকেত
ব্যবহার করেছিলেন। দুর্ভাগ্যবশত, সেখানে তিনি
$\alpha\geq\omega$ হলে $\beth_\alpha$ নেই—এমন
একটি ভুল যুক্তিও দিয়েছিলেন।}
এই সংকেতগুলি আমাদের অসীম অঙ্কবাচক সংখ্যা দেয়।

\begin{prop}
প্রত্যেক ক্রমসংখ্যা $\alpha$-র জন্য
$\aleph_\alpha$ ও $\beth_\alpha$ অঙ্কবাচক সংখ্যা।
\end{prop}

\begin{proof}
দুটি ফলই সহজ ট্রান্সফাইনাইট আবেশে পাওয়া যায়।
$\aleph_0=\beth_0=\omega$ একটি অঙ্কবাচক সংখ্যা,
\olref[cardinals][classing]{omegaisacardinal} অনুযায়ী।
$\aleph_\alpha$ ও $\beth_\alpha$ দুটিই অঙ্কবাচক সংখ্যা
ধরলে, $\aleph_{\alpha+1}$ ও $\beth_{\alpha+1}$-ও
সংজ্ঞাতেই অঙ্কবাচক সংখ্যা। আর
\olref[cardinals][classing]{unioncardinalscardinal}
অনুযায়ী অঙ্কবাচক সংখ্যার সেটের সংযোগও অঙ্কবাচক সংখ্যা।
\end{proof}
\noindent
উপরন্তু প্রত্যেক অসীম অঙ্কবাচক সংখ্যা কোনো না কোনো
$\aleph$:

\begin{prop}
$\cardfont{a}$ অসীম অঙ্কবাচক সংখ্যা হলে অদ্বিতীয়
কোনো $\gamma$-র জন্য
$\cardfont{a}=\aleph_\gamma$।
\end{prop}

\begin{proof}
অসীম অঙ্কবাচক সংখ্যার উপর ট্রান্সফাইনাইট আবেশ করি।
প্রথম ক্ষেত্রটি আগের প্রস্তাবেই প্রতিষ্ঠিত। আবেশে ধরি,
% BN-SRC-460: the predecessor induction domain excludes finite cardinals.
অসীম $\cardfont{b}<\cardfont{a}$ হলে
$\cardfont{b}=\aleph_{\gamma_\cardfont{b}}$।
কোনো অসীম $\cardfont{b}$-র জন্য
$\cardfont{a}=\cardsucc{\cardfont{b}}$ হলে
$\cardfont{a}=\cardsucc{(\aleph_{\gamma_\cardfont{b}})}=
\aleph_{\gamma_\cardfont{b}+1}$।
আর $\cardfont{a}$ কোনো অঙ্কবাচক সংখ্যার উত্তরসূরি
না হলে, অঙ্কবাচক সংখ্যাগুলি ক্রমসংখ্যা বলে
$\cardfont{a}=\bigcup_{\omega\leq\cardfont{b}<\cardfont{a}}\cardfont{b}=
\bigcup_{\omega\leq\cardfont{b}<\cardfont{a}}{\aleph_{\gamma_\cardfont{b}}}$।
তাই $\gamma=
\bigcup_{\omega\leq\cardfont{b}<\cardfont{a}}\gamma_\cardfont{b}$
নিলে $\cardfont{a}=\aleph_\gamma$।
\end{proof}

প্রত্যেক অসীম অঙ্কবাচক সংখ্যা কোনো $\aleph$ বলে
প্রশ্ন আসে: প্রত্যেকটি কি কোনো $\beth$-ও? যদি
\emph{তাই} হতো, তবে অসীম অঙ্কবাচক সংখ্যাগুলি
ঘাতসেট নেওয়ার সঙ্গে সরলভাবে সম্পর্কিত থাকত।
বিশেষত নিচের কথাটি পেতাম:

\begin{defish}
\emph{সাধারণীকৃত কন্টিনিউয়াম অনুমান} (GCH)।
প্রত্যেক $\alpha$-র জন্য $\aleph_\alpha=\beth_\alpha$।
\end{defish}

GCH সত্য হলে অঙ্কবাচক সংখ্যার ঘাতে বেশ কিছু
সরলীকরণও করা যেত। বিশেষত অসীম $\cardfont{a}$ এবং
% BN-SRC-461: exponent zero is excluded from the claimed lower bound.
অশূন্য $\cardfont{b}<\cardfont{a}$-র জন্য
$\cardexpo{\cardfont{a}}{\cardfont{b}}$-র মান
$\cardfont{a}\leq\cardexpo{\cardfont{a}}{\cardfont{b}}
\leq\cardsucc{\cardfont{a}}$-এর মধ্যে থাকত।
তারপর কোন শর্তে দুটি সম্ভাবনার কোনটি ঘটে,
অর্থাৎ $\cardfont{a}=\cardexpo{\cardfont{a}}{\cardfont{b}}$
নাকি $\cardexpo{\cardfont{a}}{\cardfont{b}}=
\cardsucc{\cardfont{a}}$, তাও নির্দিষ্টভাবে বলা যেত।
\footnote{শর্তটি \emph{সহশেষতা} দিয়ে নির্ধারিত হয়।}

কিন্তু GCH একটি \emph{অনুমান}, \emph{উপপাদ্য} নয়।
\citet{Godel1938} প্রমাণ করেন, $\ZFC$ সংগতি-সম্মত
হলে $\ZFC+\text{GCH}$-ও সংগতি-সম্মত। পরে জানা যায়,
$\ZFC$-র সঙ্গে $\lnot$GCH-ও যোগ করা যায়।
GCH-এর সবচেয়ে সরল অ-তুচ্ছ \emph{দৃষ্টান্তটি} ধরো:

\begin{defish}
\emph{কন্টিনিউয়াম অনুমান} (CH)। $\aleph_1=\beth_1$।
\end{defish}

\citet{Cohen1963} প্রমাণ করেন, $\ZFC$ সংগতি-সম্মত
হলে $\ZFC+\lnot\text{CH}$-ও সংগতি-সম্মত।
অতএব কন্টিনিউয়াম অনুমান $\ZFC$ থেকে স্বতন্ত্র।

``কন্টিনিউয়াম'' বাস্তব সংখ্যারেখা $\Real$-এর
আরেক নাম বলেই অনুমানটির এই নাম।
\olref[opps]{continuumis2aleph0} অনুযায়ী
$\card{\Real}=\beth_1$। তাই কন্টিনিউয়াম অনুমান
বলে, স্বাভাবিক সংখ্যার অঙ্কবাচকতা
$\aleph_0=\beth_0$ এবং কন্টিনিউয়ামের অঙ্কবাচকতা
$\beth_1$-এর মাঝখানে আর কোনো অঙ্কবাচক সংখ্যা নেই।

$\ZFC$ থেকে (G)CH-এর \emph{স্বতন্ত্রতা} জানা গেলে
তাদের \emph{সত্যতা} সম্পর্কে কী বলব? এ নিয়ে
অনেক কথাই বলা যায়; সেটতত্ত্বে সাম্প্রতিক বহু
ফলপ্রসূ গবেষণা এই প্রশ্নগুলির দিকে গেছে। তবে
দুটি সংক্ষিপ্ত কথা বিশেষভাবে মনে রাখা দরকার।

প্রথমত, আনুষ্ঠানিক স্বতন্ত্রতার ফল থেকে (G)CH বা
CH-এর সত্যমান \emph{অনির্ধারিত}, এমন সিদ্ধান্ত
\emph{সঙ্গে সঙ্গে} আসে না। হয়তো আমাদের কাছে
স্বাভাবিক মনে হবে এমন আরও কিছু স্বতঃসিদ্ধ যোগ
করতে হবে, এবং সেগুলি প্রশ্নটির একদিকে বা অন্যদিকে
মীমাংসা করবে। গ্যোডেল নিজে এই উত্তরই যথার্থ
বলে মনে করেছিলেন।

দ্বিতীয়ত, $\ZFC$ থেকে CH-এর স্বতন্ত্রতা অবশ্যই
\emph{চমকপ্রদ}, কিন্তু আক্ষরিক অর্থে \emph{অবিশ্বাস্য}
নয়। মূল কথা হলো, $\ZFC$ যতটুকু বলে, তাতে
অঙ্কবাচক সংখ্যা থেকে তার উত্তরসূরিতে পৌঁছানোর
পথ ঘাতসেট নেওয়ার মতো মোটাদাগের ক্রিয়ার চেয়ে
সূক্ষ্ম হতে পারে।
 %The operation of taking powersets moves us from one stage of the
 %hierarchy to its successor stage (from $V_{\alpha}$ to
 %$V_{\alpha+1}$).

এই দুটি পর্যবেক্ষণের পরে আরও জানতে চাইলে,
এই বিতর্কে অংশ নেওয়া দার্শনিক ও গণিতবিদদের
লেখার দিকেই যেতে হবে।\footnote{শুরুতে
\citet[\S15.6]{Potter2004} পড়তে পারো।}

\end{document}
